\documentclass[12pt,a4paper]{amsart}

\input{preamble}

\title[Example note]{Example note document}
\author[Author]{Author}
\date{\today}

\showdraftnotestrue
\showcommentstrue

\newauthorcomment{alice}{Alice}{blue}
\newauthorcomment{bob}{Bob}{red}

\begin{document}

\maketitle

\begin{abstract}
This short note showcases the note template: theorem-like environments,
theorem-numbered lists, ordinary enumerated lists, displayed equations,
commutative diagrams, citations, URLs, and cross-references.
\end{abstract}


\section{Objects and diagrams}

\begin{definition}
  \label{definition:good-object}
  A \emph{good object} is an object \(X\) together with a morphism
  \[
    \eta_X \colon X \to F(X)
  \]
  satisfying the following conditions.
  \draftnote{A draft note here.}
  \begin{theoremitems}
    \item The assignment \(X \mapsto \eta_X\) is functorial in \(X\).
      \label{item:good-functoriality} 
    \item The morphism \(\eta_X\) is compatible with base change.
      \label{item:good-compatibility} 
  \end{theoremitems}
\end{definition}

The numbering in \Cref{item:good-functoriality,item:good-compatibility} is
derived from the surrounding theorem-like environment
\Cref{definition:good-object}. This is useful for statements with several
numbered conditions.

\alice{We should add an example here.}

\begin{example}
  \label{example:diagram}
  A basic diagrammatic situation is a commutative square
  \[
  \begin{tikzcd}
    X' \arrow[r, "g'"] \arrow[d, "f'"'] & X \arrow[d, "f"] \\
    Y' \arrow[r, "g"']                  & Y .
  \end{tikzcd}
  \]
  The command \texttt{tikz-cd} is available in this template for diagrams of this
  kind.
\end{example}

\bob{I've added the example above.}

\begin{lemma}
  \label{lemma:base-change}
  In the situation of \Cref{example:diagram}, suppose that the square is
  cartesian. Then the compatibility condition in
  \Cref{item:good-compatibility} can be expressed by the identity
  \[
    g^* f_* \simeq f'_* {g'}^* .
  \]
\end{lemma}

\begin{proof}
  This is only a template proof. The displayed formula illustrates a typical
  mathematical display, while the reference to
  \Cref{item:good-compatibility} illustrates a theorem-numbered item reference.
\end{proof}

\section{References and ordinary lists}

The template uses biblatex. For example, one may cite
Hartshorne's textbook \cite{har77} or the Stacks Project \cite{stacks}.

Ordinary lists remain ordinary. For instance:
\begin{enumerate}[label=(\alph*), ref=(\alph*)]
  \item This is a first ordinary item.
    \label{item:ordinary-first} 
  \item This is a second ordinary item.
    \label{item:ordinary-second} 
\end{enumerate}
Thus \Cref{item:ordinary-second} is an ordinary item reference, whereas
\Cref{item:good-compatibility} is a theorem-numbered item reference.

\begin{proposition}
  \label{proposition:main}
  Let \(X\) be a good object. If the square in \Cref{example:diagram} is
  cartesian, then \(X'\) inherits the corresponding compatibility property.
\end{proposition}

\begin{proof}
  The assertion follows formally from \Cref{lemma:base-change} and
  \Cref{item:good-compatibility}.
  \bob{Maybe we should add more details here.}
\end{proof}

\begin{remark}
  \label{remark:links}
  Hyperlinks can be written using \verb|\href|. For example, the template
  repository could link to
  \href{https://www.ctan.org/pkg/tikz-cd}{the tikz-cd package page}.
  URLs inherit the hyperlink colour configured in the common preamble.
\end{remark}

\printbibliography

\end{document}
